
\documentclass{article}
\usepackage{colortbl}
\usepackage{makecell}
\usepackage{multirow}
\usepackage{supertabular}

\begin{document}

\newcounter{utterance}

\centering \large Interaction Transcript for game `wordle\_withclue', experiment `medium\_frequency\_words\_clue\_no\_critic', episode 19 with Qwen3.5{-}4B.
\vspace{24pt}

{ \footnotesize  \setcounter{utterance}{1}
\setlength{\tabcolsep}{0pt}
\begin{supertabular}{c@{$\;$}|p{.15\linewidth}@{}p{.15\linewidth}p{.15\linewidth}p{.15\linewidth}p{.15\linewidth}p{.15\linewidth}}
    \# & \multicolumn{2}{c}{Player} && \multicolumn{2}{c}{Game Master} \\
    \hline

    \theutterance \stepcounter{utterance}  
    & & & \multicolumn{4}{p{0.6\linewidth}}{
        \cellcolor[rgb]{0.9,0.9,0.9}{
            \makecell[{{p{\linewidth}}}]{
                \texttt{\tiny{[P1$\langle$GM]}}
                \texttt{You are a language wizard who likes to guess words by using the given rules.} \\
\\ 
\texttt{Welcome to Wordle! You have six attempts to guess the target word, a valid English word of five lowercase letters (a{-}z). Please use the tags "explanation:" and "guess:" to provide a concise explanation for each guess.} \\
\\ 
\texttt{To help you make an informed guess, you will receive a clue for the word, such as} \\
\texttt{clue: snowy white.} \\
\\ 
\texttt{Here is an example guess based on the clue:} \\
\texttt{explanation: In the fairy tail Snow White, the girl is killed because she eats a poisoned apple. And the word apple has 5 letters.} \\
\texttt{guess: apple} \\
\\ 
\texttt{After each guess, your answer will be validated, and you will receive feedback indicating which letters are correct (green), which letters are correct but in the wrong position (yellow), and which letters are incorrect (red). This feedback can be useful in determining which letters to include or exclude in your next guess.} \\
\\ 
\texttt{For example, the feedback for "apple" might be:} \\
\texttt{guess\_feedback: a<yellow> p<yellow> p<green> l<yellow> e<red>} \\
\\ 
\texttt{The explanation should contain details about how both the guess\_feedback and clue work together in tandem to formulate a new guess.} \\
\\ 
\texttt{Do not generate any other text.} \\
\\ 
\texttt{Let's begin with your first guess.} \\
\\ 
\texttt{clue: tend the hearth} \\
            }
        }
    }
     \\ \\

    \theutterance \stepcounter{utterance}  
    & \multicolumn{4}{p{0.6\linewidth}}{
        \cellcolor[rgb]{0.561,0.737,0.561}{
            \makecell[{{p{\linewidth}}}]{
                \texttt{\tiny{[P1$\rangle$GM]}}
                \texttt{explanation: The clue "tend the hearth" suggests a word related to caring for a fire or home, often involving a specific tool or action. The word "stove" is a common five{-}letter word associated with heating, but "stove" doesn't perfectly fit the poetic "hearth" imagery as well as "coals" or "embers". However, "embers" fits the length and the context of a hearth perfectly. Let's start with a very common five{-}letter word related to fire or heat to establish a baseline. "fire" is too short. "burn" is too short. "ash" is too short. "wood" is too short. "coal" is too short. "stove" is a strong candidate. Let's try "stove".} \\
\texttt{guess: stove} \\
            }
        }
    }
    & & \\ \\

    \theutterance \stepcounter{utterance}  
    & & & \multicolumn{2}{p{0.3\linewidth}}{
        \cellcolor[rgb]{0.9,0.9,0.9}{
            \makecell[{{p{\linewidth}}}]{
                \texttt{\tiny{[GM$|$GM]}}
                \texttt{attempts = 1} \\
\texttt{target\_word = stoke} \\
\texttt{target\_word\_clue = tend the hearth} \\
\texttt{guess = stove} \\
\texttt{guess\_feedback = s<green> t<green> o<green> v<red> e<green>} \\
            }
        }
    }
    & & \\ \\

    \theutterance \stepcounter{utterance}  
    & & & \multicolumn{4}{p{0.6\linewidth}}{
        \cellcolor[rgb]{0.9,0.9,0.9}{
            \makecell[{{p{\linewidth}}}]{
                \texttt{\tiny{[P1$\langle$GM]}}
                \texttt{guess\_feedback: s<green> t<green> o<green> v<red> e<green>} \\
\\ 
\texttt{Provide your response only in this format.} \\
\texttt{explanation: details} \\
\texttt{guess: word} \\
            }
        }
    }
     \\ \\

    \theutterance \stepcounter{utterance}  
    & \multicolumn{4}{p{0.6\linewidth}}{
        \cellcolor[rgb]{0.561,0.737,0.561}{
            \makecell[{{p{\linewidth}}}]{
                \texttt{\tiny{[P1$\rangle$GM]}}
                \texttt{explanation: The feedback s<green> t<green> o<green> v<red> e<green> confirms that the word contains the letters s, t, o, and e in those specific positions, while 'v' is incorrect. This means the word is "stoe" followed by a letter that isn't 'v'. Given the clue "tend the hearth," which often implies keeping a fire going, the word is likely "stoep" (a Dutch word) or perhaps I misidentified the target. Wait, looking at the pattern s{-}t{-}o{-}e{-}\_, the only common English word fitting this pattern related to a hearth or fire is actually "stove" but 'v' was red? That's contradictory. Let me re{-}evaluate. If s, t, o, e are green in positions 1, 2, 3, 5, the word is s{-}t{-}o{-}?{-}e. The clue "tend the hearth" might refer to "stoke". Let's check "stoke": s (green), t (green), o (green), k (unknown), e (green). This fits perfectly. The previous guess "stove" had 'v' red, so the 4th letter is not 'v'. "Stoke" fits the clue of tending a fire (stoking the fire) and matches the known letters.} \\
\texttt{guess: stoke} \\
            }
        }
    }
    & & \\ \\

    \theutterance \stepcounter{utterance}  
    & & & \multicolumn{2}{p{0.3\linewidth}}{
        \cellcolor[rgb]{0.9,0.9,0.9}{
            \makecell[{{p{\linewidth}}}]{
                \texttt{\tiny{[GM$|$GM]}}
                \texttt{attempts = 2} \\
\texttt{target\_word = stoke} \\
\texttt{target\_word\_clue = tend the hearth} \\
\texttt{guess = stoke} \\
\texttt{guess\_feedback = s<green> t<green> o<green> k<green> e<green>} \\
            }
        }
    }
    & & \\ \\

    \theutterance \stepcounter{utterance}  
    & & & \multicolumn{2}{p{0.3\linewidth}}{
        \cellcolor[rgb]{0.9,0.9,0.9}{
            \makecell[{{p{\linewidth}}}]{
                \texttt{\tiny{[GM$|$GM]}}
                \texttt{game\_result = WIN} \\
            }
        }
    }
    & & \\ \\

\end{supertabular}
}

\end{document}
